% ---------------------------------------------------------------------
% Objetivos y competencias
% ---------------------------------------------------------------------

\chapter{Objetivos y competencias}
\label{cap:objetivos-competencias}

\begin{Resumen}
En este capítulo se definen el objetivo general, los objetivos específicos y la relación del trabajo con las competencias de la titulación. También se delimitan las restricciones técnicas que condicionan el desarrollo del sistema.
\end{Resumen}

\section{Objetivo general}
\label{sec:objetivo-general}

El objetivo general del Trabajo de Fin de Grado es diseñar, ensamblar e integrar un dron personalizado que sirva como banco de pruebas para navegación autónoma basada en visión, permitiendo capturar imágenes y telemetría en vuelo real y evaluar posteriormente un algoritmo de Geolocalización Visual Cruzada.

\section{Objetivos específicos}
\label{sec:objetivos-especificos}

Para alcanzar el objetivo general se plantean los siguientes objetivos específicos:

\begin{enumerate}
\item Seleccionar e integrar los componentes mecánicos, eléctricos y electrónicos necesarios para construir una plataforma aérea funcional.
\item Configurar el controlador de vuelo, el módulo GPS con brújula y los parámetros básicos de seguridad y estabilidad.
\item Integrar una Raspberry Pi con una cámara USB de obturador global para capturar imágenes durante el vuelo.
\item Establecer un enlace de telemetría basado en MAVLink que permita registrar actitud, rumbo, altitud y posición durante la misión.
\item Preparar un flujo de trabajo para ejecutar misiones mediante estación de tierra y mantener acceso remoto mediante Tailscale y sesiones persistentes con \texttt{tmux}.
\item Construir una base de datos georreferenciada a partir de imágenes satelitales de la zona de vuelo.
\item Evaluar un pipeline de CVGL basado en recuperación global y emparejamiento local sobre las imágenes reales capturadas.
\item Analizar las limitaciones de ejecutar el procesamiento visual a bordo y proponer mejoras futuras para acercar el sistema al tiempo real.
\end{enumerate}

\section{Restricciones de diseño}
\label{sec:restricciones-diseno}

El diseño del sistema está condicionado por las restricciones propias de una plataforma aérea ligera. Las más relevantes son:

\begin{itemize}
\item \textbf{Peso}: todos los elementos embarcados reducen la autonomía y afectan a la estabilidad del dron.
\item \textbf{Energía}: la batería principal debe alimentar el sistema de propulsión y, mediante reguladores adecuados, la electrónica auxiliar.
\item \textbf{Cómputo}: la Raspberry Pi permite captura y comunicaciones, pero tiene capacidad limitada para ejecutar modelos de visión pesados.
\item \textbf{Vibraciones}: la cámara y la electrónica deben montarse minimizando vibraciones para no degradar la calidad de imagen.
\item \textbf{Comunicaciones}: el enlace remoto debe ser robusto frente a desconexiones puntuales durante la misión.
\item \textbf{Seguridad}: las pruebas deben realizarse de forma incremental, validando primero el sistema en tierra y después en vuelos controlados.
\end{itemize}

\section{Competencias relacionadas}
\label{sec:competencias-relacionadas}

El trabajo se relaciona directamente con varias competencias de la titulación. La \autoref{tab:competencias} resume la contribución prevista de cada una.

\begin{table}[htbp]
\centering
\caption{Relación entre competencias y actividades del proyecto.}
\label{tab:competencias}
\begin{tabularx}{\textwidth}{>{\bfseries}p{0.15\textwidth}X}
\toprule
Competencia & Aplicación en el TFG \\
\midrule
CE10 & Integración de una computadora de borde, cámara, telemetría y controlador de vuelo en un sistema empotrado embarcado. \\
CE14 & Uso de imágenes capturadas por una cámara para estimar localización y evaluar correspondencias visuales con mapas de referencia. \\
CE26 & Aplicación de descriptores visuales y técnicas de aprendizaje automático para reconocimiento visual de lugares. \\
CE34 & Estudio de técnicas de estimación de localización y soporte a la navegación de un robot móvil aéreo. \\
CE42 & Uso de herramientas y sistemas específicos de robótica, estaciones de tierra y protocolos de comunicación para integrar dispositivos externos. \\
\bottomrule
\end{tabularx}
\end{table}

\section{Criterios de éxito}
\label{sec:criterios-exito}

Se considerará que la integración básica del sistema es satisfactoria si el dron puede realizar vuelos controlados, capturar imágenes con marcas temporales, registrar telemetría asociada y generar un conjunto de datos utilizable por el pipeline de CVGL. En una segunda fase, el sistema se evaluará mediante métricas de error de localización, número de correspondencias geométricamente válidas y tiempo de procesamiento por imagen.

% ---------------------------------------------------------------------